\hypertarget{appendix-ea798b1ceb35}{}
\section*{Physical and logical binary states}
\section{Physical and logical binary states}\label{physical-and-logical-binary-states}

\textbf{Curriculum depth:} Introduced · \textbf{First seen:} C001 · \textbf{Latest development:} C001

\subsection{Definition and mechanism}\label{definition-and-mechanism}

Physical ranges or stored states are classified as logical 0 or 1; noise tolerance comes from accepting ranges rather than exact values.

A threshold is a boundary used to classify a measured physical quantity. A noise margin is the distance between a valid signal level and a boundary: variation inside that margin does not change the interpreted bit. A storage element retains state; combinational logic computes an output from current inputs. Exact voltage limits depend on the device.

\subsection{Prerequisites}\label{prerequisites}

No prior curriculum concept is required; this is a starting point.

\subsection{Used by}\label{used-by}

\begin{itemize}
\tightlist
\item
  \hyperref[appendix-ddbb69998a2e]{Bit} --- uses physical and logical binary states as part of its representation or reasoning.
\end{itemize}

\subsection{Recall and next depth}\label{recall-and-next-depth}

Explain physical and logical binary states without the label, work the example above, then state the assumption whose removal would change the result.

\textbf{Pending:} electrical thresholds, noise margins, storage cells, and clocked state.

\subsection{Source and navigation}\label{source-and-navigation}

\href{Continuous\%20Notes/Computing\%20Foundations\%20-\%20From\%20Bits\%20to\%20C++\%20and\%20Systems\%20APIs\#From\%20a\%20signal\%20to\%20an\%20interpretation}{C001 teaching chapter} · \hyperref[appendix-fd090d982a15]{Part definitions and mechanisms} · \href{}{Dependency Map (02-Foundations-II, appendix-4ad47bdd269a)}
